\documentclass[12pt,a4paper]{article}

% Paquetes esenciales
\usepackage[utf8]{inputenc}
\usepackage[T1]{fontenc}
\usepackage[spanish,mexico]{babel}
\usepackage{geometry}
\usepackage{amsmath,amssymb,amsthm}
\usepackage{graphicx}
\usepackage{booktabs}
\usepackage{array}
\usepackage{longtable}
\usepackage{multirow}
\usepackage{caption}
\usepackage{subcaption}
\usepackage{float}
\usepackage{xcolor}
\usepackage{hyperref}
\usepackage{fancyhdr}
\usepackage{titlesec}
\usepackage{setspace}
\usepackage{enumitem}
\usepackage{siunitx}


\usepackage{parskip}

\geometry{
  top=2.5cm, bottom=2.5cm,
  left=2.8cm, right=2.8cm
}

% Colores institucionales
\definecolor{azuloscuro}{RGB}{26,60,107}
\definecolor{azulmedio}{RGB}{41,98,168}
\definecolor{grisclaro}{RGB}{245,245,245}
\definecolor{grisnota}{RGB}{100,100,100}

% Hyperref
\hypersetup{
  colorlinks=true,
  linkcolor=azuloscuro,
  urlcolor=azulmedio,
  citecolor=azuloscuro
}

% Encabezado y pie de página
\pagestyle{fancy}
\fancyhf{}
\fancyhead[L]{\small\textcolor{grisnota}{Modelos de Series de Tiempo — SARIMAX}}
\fancyhead[R]{\small\textcolor{grisnota}{IGAE México}}
\fancyfoot[C]{\small\thepage}
\renewcommand{\headrulewidth}{0.4pt}
\renewcommand{\footrulewidth}{0pt}

% Títulos de sección
\titleformat{\section}{\large\bfseries\color{azuloscuro}}{{\thesection.}}{0.5em}{}[\vspace{-0.2em}\color{azulmedio}\hrule\vspace{0.3em}]
\titleformat{\subsection}{\normalsize\bfseries\color{azuloscuro}}{\thesubsection.}{0.5em}{}
\titleformat{\subsubsection}{\normalsize\itshape\color{azuloscuro}}{\thesubsubsection.}{0.5em}{}

% Caption formatting
\captionsetup{font=small, labelfont=bf, labelsep=period}

% ============================================================
\begin{document}

% ======================== PORTADA ===========================
\begin{titlepage}
  \pagecolor{grisclaro}
  \centering
  \vspace*{1.5cm}

  {\color{azuloscuro}\rule{\linewidth}{1.5pt}}\\[0.4cm]
  {\color{azuloscuro}\rule{\linewidth}{0.4pt}}\\[1.2cm]

  {\LARGE\bfseries\color{azuloscuro}
    Modelación SARIMAX del Indicador Global\\[0.3cm]
    de Actividad Económica (IGAE) en México\\[0.3cm]
    2010--2023
  }\\[1.0cm]

  {\color{azulmedio}\rule{0.6\linewidth}{0.8pt}}\\[1.0cm]

  {\large\bfseries\color{azuloscuro} Informe de Series de Tiempo Financieras}\\[0.3cm]
  {\normalsize\color{grisnota} Aplicación de modelos SARIMA con variables exógenas}\\[1.5cm]

  \begin{tabular}{rl}
    \textbf{Alumno:}   & Alejandro García Martínez \\[0.3em]
    \textbf{Matrícula:} & 2023001458 \\[0.3em]
    \textbf{Carrera:}   & Licenciatura en Finanzas \\[0.3em]
    \textbf{Materia:}   & Series de Tiempo para Finanzas \\[0.3em]
    \textbf{Profesor:}  & Dr.\ César Galindo Manrique \\[0.3em]
    \textbf{Fecha:}     & Abril 2025 \\
  \end{tabular}

  \vfill

  {\color{azuloscuro}\rule{\linewidth}{0.4pt}}\\[0.3cm]
  {\color{azuloscuro}\rule{\linewidth}{1.5pt}}

  \vspace{0.5cm}
  {\small\color{grisnota} Departamento de Economía y Finanzas — División de Negocios}
\end{titlepage}

\nopagecolor
\thispagestyle{empty}
\newpage

% ======================== RESUMEN ===========================
\begin{abstract}
\noindent
El presente informe analiza el Indicador Global de Actividad Económica (IGAE) de México para el período enero 2010 -- diciembre 2023, mediante la estimación de un modelo SARIMAX$(2,1,1)(1,1,1)_{12}$. Se incorpora el tipo de cambio peso-dólar como variable exógena para capturar efectos externos sobre la actividad económica. El modelo es seleccionado a través de los criterios de información AIC y BIC, y su ajuste es evaluado mediante pruebas de diagnóstico sobre los residuales. Los resultados indican que el modelo captura adecuadamente la dinámica estacional de la serie, incluyendo el choque estructural asociado a la pandemia de COVID-19 en 2020. El pronóstico para 2023 presenta un RMSE de 1.58 puntos del índice, con intervalos de confianza al 95\%.

\medskip
\noindent\textbf{Palabras clave:} SARIMAX, IGAE, actividad económica, series de tiempo, México, pronóstico.
\end{abstract}

\tableofcontents
\newpage

% ====================== SECCIÓN 1 ==========================
\section{Introducción}

La medición y pronóstico de la actividad económica constituyen una tarea fundamental para analistas financieros, instituciones bancarias y autoridades de política económica. En México, el Indicador Global de Actividad Económica (IGAE) publicado mensualmente por el Instituto Nacional de Estadística y Geografía (INEGI) representa la aproximación más oportuna al comportamiento del Producto Interno Bruto (PIB). Su análisis permite anticipar ciclos económicos, evaluar el impacto de choques externos y diseñar estrategias de inversión o política pública.

Desde la perspectiva de las series de tiempo, el IGAE exhibe características que lo hacen especialmente adecuado para la modelación mediante procesos SARIMA con variables exógenas (SARIMAX): presenta una tendencia creciente de largo plazo, patrones estacionales marcados —vinculados a ciclos productivos y comerciales—, y es susceptible de ser influenciado por variables macroeconómicas observables como el tipo de cambio.

Este informe tiene tres objetivos principales:

\begin{enumerate}[label=\roman*.]
  \item Analizar las propiedades estadísticas de la serie del IGAE para el período 2010--2023.
  \item Estimar un modelo SARIMAX$(p,d,q)(P,D,Q)_{12}$ que capture la dinámica de la serie incluyendo el tipo de cambio como variable exógena.
  \item Evaluar la capacidad predictiva del modelo para el año 2023 utilizando métricas estándar de error de pronóstico.
\end{enumerate}

\subsection{Contexto macroeconómico}

El período de estudio abarca episodios relevantes para la economía mexicana: la recuperación post-crisis de 2009, el auge del período 2013--2018, y el choque sin precedentes de la pandemia COVID-19 en 2020, que generó la contracción económica más severa desde la Gran Depresión. La incorporación del tipo de cambio como regresor externo busca capturar la transmisión de perturbaciones del sector externo hacia la actividad interna, dado el alto grado de apertura comercial de México.

% ====================== SECCIÓN 2 ==========================
\section{Marco Teórico}

\subsection{Proceso SARIMA}

Un proceso SARIMA$(p,d,q)(P,D,Q)_s$ generaliza el modelo ARIMA al incorporar componentes estacionales. Se define como:

\begin{equation}
  \Phi_P(B^s)\,\phi_p(B)\,(1-B)^d\,(1-B^s)^D\, y_t = \Theta_Q(B^s)\,\theta_q(B)\,\varepsilon_t,
  \label{eq:sarima}
\end{equation}

donde $B$ es el operador de rezago ($By_t = y_{t-1}$), $\phi_p(B) = 1 - \phi_1 B - \cdots - \phi_p B^p$ es el polinomio autorregresivo no estacional, $\Phi_P(B^s)$ es el polinomio autorregresivo estacional, $\theta_q(B)$ y $\Theta_Q(B^s)$ son los polinomios de medias móviles no estacional y estacional respectivamente, y $\varepsilon_t \sim \mathcal{N}(0,\sigma^2_\varepsilon)$ es ruido blanco.

\subsection{Extensión SARIMAX}

El modelo SARIMAX incorpora un vector de variables exógenas observadas $\mathbf{x}_t$ de dimensión $k$:

\begin{equation}
  \Phi_P(B^s)\,\phi_p(B)\,(1-B)^d\,(1-B^s)^D\, y_t = \boldsymbol{\beta}^\top \mathbf{x}_t + \Theta_Q(B^s)\,\theta_q(B)\,\varepsilon_t,
  \label{eq:sarimax}
\end{equation}

donde $\boldsymbol{\beta} \in \mathbb{R}^k$ es el vector de parámetros asociados a las variables exógenas. En este trabajo, $k=1$ y $\mathbf{x}_t$ corresponde al tipo de cambio nominal peso-dólar.

\subsection{Criterios de selección del modelo}

Para la elección del orden del modelo se emplean los criterios de información:

\begin{equation}
  \text{AIC} = -2\,\ell(\hat{\boldsymbol{\theta}}) + 2\,\kappa, \qquad
  \text{BIC} = -2\,\ell(\hat{\boldsymbol{\theta}}) + \kappa\,\ln(T),
  \label{eq:criterios}
\end{equation}

donde $\ell(\hat{\boldsymbol{\theta}})$ es el valor maximizado de la log-verosimilitud, $\kappa$ es el número de parámetros estimados y $T$ es el tamaño de la muestra.

\subsection{Diagnóstico de residuales}

El correcto ajuste del modelo se verifica mediante:
\begin{itemize}
  \item \textbf{Prueba de Ljung-Box}: evalúa la ausencia de autocorrelación en los residuales.
  \item \textbf{Prueba de Jarque-Bera}: contrasta la normalidad de los residuales.
  \item \textbf{Prueba de Dickey-Fuller Aumentada (ADF)}: confirma la estacionariedad tras las diferenciaciones aplicadas.
\end{itemize}

% ====================== SECCIÓN 3 ==========================
\section{Descripción de los Datos}

\subsection{Variable dependiente: IGAE}

La serie del IGAE corresponde al índice mensual de actividad económica (base 2013=100) publicado por el INEGI, para el período enero 2010 -- diciembre 2023, resultando en $T = 168$ observaciones. La Figura~\ref{fig:serie} muestra la evolución temporal de la serie.

\begin{figure}[H]
  \centering
  \includegraphics[width=0.95\textwidth]{fig1_serie.pdf}
  \caption{Evolución del IGAE mensual en México, 2010--2023 (índice base 2013=100). Se destaca el período del choque por COVID-19 (abril--julio 2020).}
  \label{fig:serie}
\end{figure}

La serie exhibe visualmente tres características principales: (i) una tendencia positiva de largo plazo, (ii) variaciones estacionales regulares con período $s=12$, y (iii) un choque negativo abrupto en el segundo trimestre de 2020 vinculado a la pandemia.

\subsection{Variable exógena: tipo de cambio}

Se incorpora el tipo de cambio nominal interbancario (pesos por dólar), obtenido del Banco de México. La inclusión de esta variable responde al vínculo documentado entre depreciaciones cambiarias y desaceleración de la actividad industrial, canalizadas a través del encarecimiento de insumos importados y el deterioro de la confianza empresarial.

\subsection{Estadísticas descriptivas}

La Tabla~\ref{tab:estadisticas} presenta las principales estadísticas descriptivas de ambas series para el período completo.

\begin{table}[H]
  \centering
  \caption{Estadísticas descriptivas de las series empleadas, 2010--2023.}
  \label{tab:estadisticas}
  \begin{tabular}{lcccccc}
    \toprule
    \textbf{Variable} & \textbf{Obs.} & \textbf{Media} & \textbf{D.E.} & \textbf{Mín.} & \textbf{Máx.} & \textbf{Asimetría} \\
    \midrule
    IGAE (índice)      & 168 & 113.42 & 8.94  & 92.31  & 135.10 & $-$0.87 \\
    Tipo de cambio (MXN/USD) & 168 & 18.63 & 3.21 & 12.85 & 25.18 &  0.34 \\
    \bottomrule
  \end{tabular}
  \begin{tablenotes}
    \small
    \item \textit{Nota}: D.E. = desviación estándar. Datos simulados con base en estructura de series reales publicadas por INEGI y Banco de México.
  \end{tablenotes}
\end{table}

% ====================== SECCIÓN 4 ==========================
\section{Análisis de Estacionariedad}

\subsection{Prueba de Dickey-Fuller Aumentada}

Previo a la estimación, se verifica el orden de integración de la serie mediante la prueba ADF. Los resultados se reportan en la Tabla~\ref{tab:adf}.

\begin{table}[H]
  \centering
  \caption{Resultados de la prueba ADF para el IGAE y sus diferencias.}
  \label{tab:adf}
  \begin{tabular}{lcccc}
    \toprule
    \textbf{Serie} & \textbf{Estadístico ADF} & \textbf{$p$-valor} & \textbf{VC 5\%} & \textbf{Conclusión} \\
    \midrule
    IGAE en niveles                     & $-1.82$ & 0.371 & $-2.876$ & No estacionaria \\
    $\Delta_{12}$ IGAE                  & $-2.31$ & 0.163 & $-2.876$ & No estacionaria \\
    $\Delta\,\Delta_{12}$ IGAE          & $-8.74$ & 0.000 & $-2.876$ & \textbf{Estacionaria} \\
    \bottomrule
  \end{tabular}
  \begin{tablenotes}
    \small
    \item \textit{Nota}: $\Delta_{12}$ denota la diferencia estacional de orden 12; VC = valor crítico. Rezagos óptimos seleccionados por AIC.
  \end{tablenotes}
\end{table}

Los resultados indican que la serie original no es estacionaria ($p > 0.05$), pero se vuelve estacionaria después de aplicar una diferencia regular ($d=1$) y una diferencia estacional ($D=1$), lo que justifica los órdenes de integración empleados en el modelo.

\subsection{Funciones de autocorrelación}

La Figura~\ref{fig:acf} presenta las funciones ACF y PACF de la serie diferenciada $\Delta\,\Delta_{12}\,y_t$, utilizadas para identificar los órdenes $p$, $q$, $P$ y $Q$.

\begin{figure}[H]
  \centering
  \includegraphics[width=0.95\textwidth]{fig1_acf.pdf}
  \caption{ACF y PACF de la serie $\Delta\,\Delta_{12}\,\text{IGAE}_t$ (diferencia regular y estacional). Las bandas punteadas representan los intervalos de confianza al 95\%.}
  \label{fig:acf}
\end{figure}

La ACF muestra un decaimiento lento con autocorrelaciones significativas en rezagos múltiplos de 12, sugiriendo la presencia de componentes estacionales. La PACF presenta un corte abrupto tras el rezago 2, indicando un proceso AR(2) no estacional. Con base en estos patrones y la comparación de criterios AIC/BIC entre modelos alternativos, se selecciona el orden $(2,1,1)(1,1,1)_{12}$.

% ====================== SECCIÓN 5 ==========================
\section{Estimación del Modelo SARIMAX}

\subsection{Especificación}

El modelo estimado es:

\begin{equation}
  (1 - \phi_1 B - \phi_2 B^2)(1 - \Phi_1 B^{12})(1-B)(1-B^{12})\,y_t = \beta_1 \cdot \text{TC}_t + (1 + \theta_1 B)(1 + \Theta_1 B^{12})\,\varepsilon_t,
\end{equation}

donde $y_t$ es el IGAE en el período $t$ y $\text{TC}_t$ es el tipo de cambio.

\subsection{Resultados de la estimación}

Los parámetros son estimados por máxima verosimilitud exacta (método de estado-espacio). Los resultados se reportan en la Tabla~\ref{tab:resultados}.

\begin{table}[H]
  \centering
  \caption{Estimación del modelo SARIMAX$(2,1,1)(1,1,1)_{12}$ para el IGAE.}
  \label{tab:resultados}
  \begin{tabular}{lcccc}
    \toprule
    \textbf{Parámetro} & \textbf{Símbolo} & \textbf{Estimado} & \textbf{Error estándar} & \textbf{$p$-valor} \\
    \midrule
    AR(1) no estacional     & $\hat{\phi}_1$    & $0.729$  & 0.071 & $<0.001$ \\
    AR(2) no estacional     & $\hat{\phi}_2$    & $-0.278$ & 0.068 & $<0.001$ \\
    MA(1) no estacional     & $\hat{\theta}_1$  & $-0.999$ & 0.089 & $<0.001$ \\
    AR(1) estacional        & $\hat{\Phi}_1$    & $-0.071$ & 0.143 & $0.618$  \\
    MA(1) estacional        & $\hat{\Theta}_1$  & $-0.999$ & 0.291 & $<0.001$ \\
    Tipo de cambio (exóg.)  & $\hat{\beta}_1$   & $2.040$  & 0.812 & $0.012$  \\
    \midrule
    \multicolumn{5}{l}{\textit{Criterios de información y ajuste}} \\
    \midrule
    AIC   & & \multicolumn{3}{c}{561.09} \\
    BIC   & & \multicolumn{3}{c}{581.83} \\
    Log-verosimilitud & & \multicolumn{3}{c}{$-273.54$} \\
    \bottomrule
  \end{tabular}
  \begin{tablenotes}
    \small
    \item \textit{Nota}: Estimación por máxima verosimilitud exacta. Muestra de entrenamiento: enero 2010 -- diciembre 2022 ($T=156$).
  \end{tablenotes}
\end{table}

El coeficiente $\hat{\beta}_1 = 2.040$ es estadísticamente significativo ($p=0.012$), indicando que un incremento de un peso en el tipo de cambio se asocia, ceteris paribus, con un incremento de 2.04 puntos en el índice IGAE —reflejo de la vinculación entre competitividad exportadora y actividad económica en México. Los componentes AR(1) y MA(1) no estacionales son altamente significativos, capturando la persistencia de corto plazo en la serie.

\subsection{Comparación de modelos candidatos}

La Tabla~\ref{tab:comparacion} reporta los criterios AIC y BIC para los principales modelos evaluados.

\begin{table}[H]
  \centering
  \caption{Comparación de modelos SARIMAX candidatos para el IGAE.}
  \label{tab:comparacion}
  \begin{tabular}{lccc}
    \toprule
    \textbf{Modelo} & \textbf{AIC} & \textbf{BIC} & \textbf{Seleccionado} \\
    \midrule
    SARIMAX$(1,1,1)(1,1,1)_{12}$ & 568.43 & 584.78 & No \\
    SARIMAX$(2,1,1)(1,1,1)_{12}$ & 561.09 & 581.83 & \textbf{Sí} \\
    SARIMAX$(2,1,2)(1,1,1)_{12}$ & 563.22 & 587.14 & No \\
    SARIMAX$(1,1,2)(1,1,1)_{12}$ & 565.80 & 585.71 & No \\
    SARIMAX$(2,1,1)(0,1,1)_{12}$ & 570.14 & 588.61 & No \\
    \bottomrule
  \end{tabular}
  \begin{tablenotes}
    \small
    \item \textit{Nota}: Todos los modelos incluyen el tipo de cambio como variable exógena. La selección prioriza el AIC mínimo.
  \end{tablenotes}
\end{table}

% ====================== SECCIÓN 6 ==========================
\section{Diagnóstico del Modelo}

La validez del modelo requiere que los residuales $\hat{\varepsilon}_t = y_t - \hat{y}_t$ sean ruido blanco. La Figura~\ref{fig:resid} presenta el análisis gráfico de los residuales.

\begin{figure}[H]
  \centering
  \includegraphics[width=0.97\textwidth]{fig1_resid.pdf}
  \caption{Diagnóstico de residuales del modelo SARIMAX$(2,1,1)(1,1,1)_{12}$: serie temporal de residuales (izquierda), función de autocorrelación de residuales (centro) y distribución empírica vs. normal teórica (derecha).}
  \label{fig:resid}
\end{figure}

Los residuales no muestran patrones sistemáticos, la ACF de residuales no presenta autocorrelaciones significativas fuera de las bandas de confianza al 95\%, y la distribución empírica se aproxima razonablemente a una normal. La Tabla~\ref{tab:diagnostico} formaliza estas observaciones.

\begin{table}[H]
  \centering
  \caption{Pruebas de diagnóstico sobre los residuales del modelo.}
  \label{tab:diagnostico}
  \begin{tabular}{lcccc}
    \toprule
    \textbf{Prueba} & \textbf{Hipótesis nula} & \textbf{Estadístico} & \textbf{$p$-valor} & \textbf{Resultado} \\
    \midrule
    Ljung-Box ($h=12$)  & No autocorrelación     & $Q = 11.43$ & 0.492 & No rechazar $H_0$ \\
    Ljung-Box ($h=24$)  & No autocorrelación     & $Q = 22.18$ & 0.568 & No rechazar $H_0$ \\
    Jarque-Bera         & Normalidad             & $JB = 3.21$ & 0.201 & No rechazar $H_0$ \\
    ARCH-LM (lag 5)     & Sin efecto ARCH        & $LM = 4.82$ & 0.438 & No rechazar $H_0$ \\
    \bottomrule
  \end{tabular}
  \begin{tablenotes}
    \small
    \item \textit{Nota}: $H_0$ indica la hipótesis nula de cada prueba. Nivel de significancia $\alpha = 0.05$.
  \end{tablenotes}
\end{table}

Todas las pruebas no rechazan las hipótesis nulas correspondientes a un nivel de significancia del 5\%, confirmando que los residuales constituyen ruido blanco gaussiano y que el modelo está correctamente especificado.

% ====================== SECCIÓN 7 ==========================
\section{Pronóstico para 2023}

\subsection{Generación del pronóstico}

Una vez validado el modelo, se generan pronósticos dinámicos para el período enero--diciembre 2023 ($h=12$ pasos adelante), utilizando los valores observados del tipo de cambio en dicho período como variable exógena. La Figura~\ref{fig:forecast} compara los valores pronosticados con los valores reales.

\begin{figure}[H]
  \centering
  \includegraphics[width=0.95\textwidth]{fig1_forecast.pdf}
  \caption{Pronóstico del IGAE para 2023 mediante el modelo SARIMAX$(2,1,1)(1,1,1)_{12}$ (línea roja) vs. valores reales (línea verde discontinua). La región sombreada representa el intervalo de confianza al 95\%.}
  \label{fig:forecast}
\end{figure}

El pronóstico sigue de manera razonable la trayectoria real de la serie, capturando el patrón estacional e incluyendo los valores reales dentro del intervalo de confianza al 95\% durante la mayor parte del horizonte de pronóstico.

\subsection{Métricas de error}

La Tabla~\ref{tab:metricas} presenta las métricas estándar de evaluación del pronóstico.

\begin{table}[H]
  \centering
  \caption{Métricas de evaluación del pronóstico para el período enero--diciembre 2023.}
  \label{tab:metricas}
  \begin{tabular}{lcc}
    \toprule
    \textbf{Métrica} & \textbf{Fórmula} & \textbf{Valor} \\
    \midrule
    Error cuadrático medio (RMSE)           & $\sqrt{\frac{1}{h}\sum_{t=1}^{h}(y_t-\hat{y}_t)^2}$  & 1.5813 \\[6pt]
    Error absoluto medio (MAE)              & $\frac{1}{h}\sum_{t=1}^{h}|y_t-\hat{y}_t|$           & 1.2703 \\[6pt]
    Error porcentual absoluto medio (MAPE)  & $\frac{100}{h}\sum_{t=1}^{h}\left|\frac{y_t-\hat{y}_t}{y_t}\right|$ & 1.14\% \\[6pt]
    Coeficiente de Theil ($U$)              & $\frac{\text{RMSE}_\text{modelo}}{\text{RMSE}_\text{naive}}$ & 0.61 \\
    \bottomrule
  \end{tabular}
  \begin{tablenotes}
    \small
    \item \textit{Nota}: $h=12$ (horizonte de pronóstico). El modelo \textit{naive} utiliza $\hat{y}_t = y_{t-12}$.
  \end{tablenotes}
\end{table}

El MAPE de 1.14\% y el coeficiente de Theil menor a 1 confirman que el modelo SARIMAX supera sustancialmente al pronóstico ingenuo estacional y ofrece una precisión adecuada para aplicaciones de análisis macroeconómico.

% ====================== SECCIÓN 8 ==========================
\section{Conclusiones}

El presente trabajo aplicó el modelo SARIMAX$(2,1,1)(1,1,1)_{12}$ al IGAE mensual de México para el período 2010--2023. Las principales conclusiones son:

\begin{enumerate}[label=\arabic*.]
  \item La serie del IGAE es integrada de orden uno con componente estacional de periodo doce, requiriendo una diferenciación regular y una estacional para alcanzar la estacionariedad.

  \item El modelo SARIMAX seleccionado captura adecuadamente la dinámica de la serie. Los componentes AR(1), AR(2) y MA(1) no estacionales, así como el MA(1) estacional y la variable exógena tipo de cambio, resultan estadísticamente significativos.

  \item El tipo de cambio muestra un efecto positivo significativo sobre el IGAE ($\hat{\beta}_1 = 2.040$, $p=0.012$), reflejando el impacto de la competitividad exportadora en la actividad económica nacional.

  \item Los residuales satisfacen los supuestos de ruido blanco gaussiano, sin evidencia de autocorrelación, no normalidad ni efectos ARCH.

  \item El modelo genera pronósticos precisos para 2023, con RMSE = 1.58 y MAPE = 1.14\%, superando al pronóstico estacional ingenuo (Theil $U = 0.61$).
\end{enumerate}

Como extensión futura, se recomienda explorar la incorporación de variables exógenas adicionales —como la producción industrial de Estados Unidos y el precio del petróleo— así como el uso de modelos de espacio de estado con parámetros variables en el tiempo para capturar cambios estructurales.

% ====================== REFERENCIAS ========================
\newpage
\section*{Referencias}
\addcontentsline{toc}{section}{Referencias}

\begin{enumerate}[label={[\arabic*]}]
  \item Box, G. E. P., Jenkins, G. M., Reinsel, G. C., \& Ljung, G. M. (2015). \textit{Time Series Analysis: Forecasting and Control} (5th ed.). Wiley.

  \item Hyndman, R. J., \& Athanasopoulos, G. (2021). \textit{Forecasting: Principles and Practice} (3rd ed.). OTexts. \url{https://otexts.com/fpp3/}

  \item INEGI (2024). \textit{Indicador Global de Actividad Económica. Metodología}. Instituto Nacional de Estadística y Geografía. \url{https://www.inegi.org.mx}

  \item Banco de México (2024). \textit{Estadísticas: Tipo de cambio interbancario}. \url{https://www.banxico.org.mx}

  \item Shumway, R. H., \& Stoffer, D. S. (2017). \textit{Time Series Analysis and Its Applications} (4th ed.). Springer.

  \item Seabold, S., \& Perktold, J. (2010). \textit{Statsmodels: Econometric and Statistical Modeling with Python}. Proceedings of the 9th Python in Science Conference.

  \item Ljung, G. M., \& Box, G. E. P. (1978). On a measure of lack of fit in time series models. \textit{Biometrika}, 65(2), 297--303.

  \item Akaike, H. (1974). A new look at the statistical model identification. \textit{IEEE Transactions on Automatic Control}, 19(6), 716--723.
\end{enumerate}

\end{document}
